% !TeX program = xelatex
% !TeX program = xelatex
% ===========================================================================
%  tech_common.tex -- 两份技术文档共用的导言区
% ===========================================================================
\documentclass[11pt,a4paper]{ctexart}

\usepackage[top=2.4cm,bottom=2.4cm,left=2.4cm,right=2.4cm]{geometry}
\usepackage{amsmath,amssymb,amsthm}
\usepackage{graphicx}
\usepackage{booktabs}
\usepackage{longtable}
\usepackage{array}
\usepackage{enumitem}
\usepackage{xcolor}
\usepackage{listings}
\usepackage{tcolorbox}
\usepackage{caption}
\usepackage{fancyhdr}
\usepackage{titlesec}
\usepackage{hyperref}

\hypersetup{colorlinks=true,linkcolor=blue!55!black,citecolor=blue!55!black,
            urlcolor=blue!55!black,bookmarksnumbered=true}

% ---------------------------------------------------------------- 字体/版式
\setlength{\emergencystretch}{2.5em}
\setlength{\parskip}{0.35em}
\setlength{\parindent}{2em}
\linespread{1.12}

% ---------------------------------------------------------------- 代码样式
\definecolor{codebg}{RGB}{247,247,250}
\definecolor{codekw}{RGB}{0,90,160}
\definecolor{codecm}{RGB}{0,128,96}
\definecolor{codest}{RGB}{170,60,60}
\lstset{
  basicstyle=\ttfamily\footnotesize,
  backgroundcolor=\color{codebg},
  frame=single, framerule=0.4pt, rulecolor=\color{gray!45},
  breaklines=true, breakatwhitespace=false,
  numbers=left, numberstyle=\tiny\color{gray!80}, numbersep=6pt,
  keywordstyle=\color{codekw}\bfseries, commentstyle=\color{codecm}\itshape,
  stringstyle=\color{codest}, showstringspaces=false, columns=flexible,
  tabsize=2, extendedchars=true, captionpos=b,
  xleftmargin=1.6em, framexleftmargin=1.4em
}
\renewcommand{\lstlistingname}{代码}

% ---------------------------------------------------------------- 提示框
\newtcolorbox{keybox}[1][]{colback=blue!4,colframe=blue!45!black,
  boxrule=0.6pt,arc=2pt,left=5pt,right=5pt,top=4pt,bottom=4pt,
  fonttitle=\bfseries,#1}
\newtcolorbox{warnbox}[1][]{colback=orange!6,colframe=orange!55!black,
  boxrule=0.6pt,arc=2pt,left=5pt,right=5pt,top=4pt,bottom=4pt,
  fonttitle=\bfseries,#1}

% ---------------------------------------------------------------- 定理环境
\newtheorem{theorem}{定理}
\newtheorem{proposition}{命题}
\newtheorem{lemma}{引理}
\renewcommand{\proofname}{\heiti 证明}

% ---------------------------------------------------------------- 页眉页脚
\pagestyle{fancy}
\fancyhf{}
\fancyhead[L]{\small\color{gray!70}\leftmark}
\fancyhead[R]{\small\color{gray!70}\thepage}
\renewcommand{\headrulewidth}{0.3pt}
\renewcommand{\headrule}{\hbox to\headwidth{\color{gray!40}\leaders\hrule height \headrulewidth\hfill}}

\titleformat{\section}{\normalfont\heiti\Large}{\thesection}{0.6em}{}
\titleformat{\subsection}{\normalfont\heiti\large}{\thesubsection}{0.6em}{}
\titleformat{\subsubsection}{\normalfont\heiti\normalsize}{\thesubsubsection}{0.6em}{}

\captionsetup{font=small,labelfont=bf,labelsep=quad}


\title{\heiti 更聪明的算法：三条方向的攻坚结果\\[4pt]
       \large 贴边定向源的严格解法、检测次数的下界、顺路清除的否证、\\
       以及"完成度 vs 时间"的真实前沿\\[2pt]
       \normalsize —— 四池 120 例联合验证，含全部负结果}
\author{2026 高教社杯全国大学生数学建模竞赛 B 题}
\date{\today}

\begin{document}
\maketitle
\thispagestyle{fancy}

\begin{abstract}
\noindent
本轮针对三条指定方向做了实质性算法改造，并在\textbf{四个互不重叠的独立算例池}
（种子 62000+、51000+、12000+、30000+，各 30 例，共 120 例约 1560 个干扰源）上联合验证。
选型一律以\textbf{最差池}为准。

\textbf{（一）贴边定向源——解决到几何极限。} 重新刻画了问题：机器人在边界点听到源后，
沿射线走不改变方位（拿不到第二条独立方位），垂直走可能越界（落到背光侧）。
正确做法是\textbf{对横向偏移量本身做二分}——因为受光半平面必然包含听到点，
偏移趋于零时必然仍可听，所以这套二分\textbf{收敛性有保证}。
配合\textbf{末段栅格穷举清除}（间距 15 m，$15/\sqrt2=10.6$ m $<$ 20 m 清除半径，
因此只要 $\sigma$ 覆盖真值就必然命中），把四个池的漏源从
\textbf{3 个降到 1 个}，而且平均时间还略降（885.3 s $\to$ 880.7 s）。

仅剩的 1 例经几何核算属于\textbf{问题本身的极限}而非算法缺陷：该源位于 $(1767,-219)$、
辐射方向与朝外径向的夹角余弦为 0.996（几乎完全朝外），
其\textbf{域内受光区只有 0.011 km$^2$，占全域的 0.11\%，是一条 27 m 宽的月牙}。
要保证覆盖它需要在整条 11.3 km 的边界上每 25 m 布一个站位（约 450 站），不可行。

\textbf{（二）减少每站检测频道数——已做到接近下界。} 实测每站检测次数随频道逐个解算
由 20 降到 9$\sim$16。新增的"站位附近已全部认证则不再测量"规则在逻辑上严格成立
（认证扫描是精确的），但实测几乎不触发——因为认证本身已接近信息下界：
每个候选点只需 3$\sim$4 条环绕读数，而格点扫描正好提供这么多。继续压缩只能靠
\textbf{削弱认证本身}，那是用完成率换时间。

\textbf{（三）顺路清除——实测为负优化。} 关闭时完成率 1.0000、540.3 s；
开启（半径 520 m）后降到 0.9946、573.3 s。原因是清除任务自身的归航与末段图案
带来的额外移动\textbf{超过了它省下的那趟绕行}，并且打断了扫描的发现进度。

\textbf{（四）混合策略的真正答案。} 论文配置的 24 个站位里有 17 个是认证搜索补出来的，
而认证搜索把站位放到"证书最弱"处——\textbf{那正是定向源可能藏身之处}。
换句话说认证搜索兼任了定向源发现器，这是它真正的价值。
把认证搜索去掉换成"更密的纯格点扫描"，时间确实降到 533.8 s，
但最差完成率从 0.9978 掉到 0.9833（120 例中 14 例失败）。
加贴边环（6$\sim$16 站）只能把最差补到 0.9905，代价是时间回到 583 s。

\textbf{最终结论（诚实版）}：在\textbf{完成度完整}的前提下，本轮把"漏源"从 3 个压到 1 个、
时间基本持平；而\textbf{"完成度完整"与"时间大幅压缩"没有同时达成}——
这是一个真实的前沿：约 880 s 对应 0.9978，约 534 s 对应 0.9833。
本报告给出完整前沿与每一项尝试的判定，供权衡。
\end{abstract}

\tableofcontents
\newpage

% ===========================================================================
\section{验证协议：四池联合，以最差池为准}

本轮的所有结论都在四个\textbf{互不重叠}的池上取平均，并以\textbf{最差池}作为选型依据：

\begin{center}
\begin{tabular}{lll}
\toprule
池 & 种子 & 规模 \\ \midrule
A & 62000--62029 & 30 例 \\
B & 51000--51029 & 30 例 \\
C & 12000--12019 & 30 例 \\
D & 30000--30029 & 30 例 \\
\bottomrule
\end{tabular}
\end{center}

\textbf{为什么必须这样。} 本项目曾三次被单一池误导：

\begin{enumerate}[leftmargin=1.8em]
  \item 第一轮"下降 44\%"的假提升（算例池被复用）；
  \item 第二轮"50\% 相关性最差"的假异常（误差分布标准差不同，对照不公平）；
  \item 第三轮本轮开始时：$s=900$ 在调参池（30000+）上是 1.0000 / 540.3 s，
        在池 A（62000+）上却只有 0.9928；而 $s=1000$ 在池 A 上是 1.0000 / 539.3 s，
        在调参池上只有 0.9831。\textbf{两个方向都被误导过。}
\end{enumerate}

因此本轮不再使用"调参池"，\textbf{所有选型直接看四池汇总}。

% ===========================================================================
\section{方向一：贴边定向源的严格解法}

\subsection{问题的完整刻画}

前三轮已经把现象描清楚了，这里给出严格的因果链：

\begin{enumerate}[leftmargin=1.8em]
  \item 一个贴在边界、朝外辐射的定向源，其\textbf{域内受光区是一条薄月牙}；
  \item 机器人只能在月牙里的某点 $p$ 听到它一次；
  \item 只有一条方位 $\Rightarrow$ 定位区域\textbf{无界}（实测 \texttt{bounded=False}）
        $\Rightarrow$ 不生成清除任务；
  \item 要拿到第二条方位就必须换观测点，而：
        \begin{itemize}
          \item \textbf{沿射线走}：方位不变（源就在这条射线上），拿不到独立的第二条方位；
          \item \textbf{垂直走}：可能越出受光半平面，返回 no\_signal。
        \end{itemize}
\end{enumerate}

前两轮的失败尝试正好落在这两个陷阱里：沿射线 550 m 单跳越过了源落到背光侧；
沿射线 200 m 爬行虽然拿到了第二条方位，但两条方位只差 $1.3^{\circ}$，
交会区域仍然无界（$\sigma=670$ m）。

\subsection{正确构造：对横向偏移量二分}

关键几何事实：

\begin{keybox}[title=偏移二分必然收敛]
机器人在 $p$ 听到源，说明 $p$ 落在源的受光半平面 $H$ 内。
取 $p$ 处垂直于视线的方向，偏移 $\delta$ 后得到 $p'$。
当 $\delta\to 0$ 时 $p'\to p\in H$，因此\textbf{存在 $\delta_0>0$，
使得所有 $\delta<\delta_0$ 的偏移点仍在 $H$ 内、仍能听到源}。
于是按 $\delta=400,250,150,90,55,33,20$ m 依次尝试（两侧都试），
\textbf{必在有限步内成功}——这不是启发式，而是由半平面的开性保证的收敛构造。
\end{keybox}

代码实现为 \texttt{relocate} 任务（\texttt{\_relocate\_point} / \texttt{\_relocate\_advance}）。
若全部偏移档位都失败，则沿射线前爬 \texttt{relocate\_crawl\_step} $=200$ m
再重新从大偏移试起——此时横向偏移的绝对尺度自动变小，等价于把二分继续下去。

\begin{warnbox}[title=实现中的一个真实 bug（值得记录）]
初版把状态机的推进写在 \texttt{\_relocate\_point()} 里，而这个函数是在\textbf{构建任务列表}
时被调用的；\texttt{\_build\_tasks()} 又在每次重规划（即每个动作之后）被调用。
结果是：状态机在\textbf{没有任何候选点被实际访问}的情况下就跑完了，
6 个优质候选点（离源 40$\sim$90 m 且受光）一个都没去。
现象是"机制明明生成了正确点位，机器人却始终停在原地"。
修法：查询与推进分离——\texttt{\_relocate\_point()} 只读取状态，
状态在\textbf{任务执行完}之后由 \texttt{\_relocate\_advance()} 推进。
\end{warnbox}

\subsection{第二件工具：末段栅格穷举清除}

诊断（池 C / seed 12010 / ch12）：源在 $r=1799$、定向，机器人已把它\textbf{定位到
$\sigma=42$ m}，但 14 次清除全部失败并用尽了硬上限，源被漏掉。
原因是末段图案取半径 $0.55\sigma\approx 23$ m、两圈 23 m 与 46 m，
图案点之间的空隙可能整片落在 20 m 清除半径之外。

\begin{keybox}[title=有保证的兜底]
当估计已足够好（$\sigma$ 不超过定位门槛）而常规末段图案连续失败后，
在 $\sigma$ 圆盘上做间距 $g=15$ m 的方形栅格穷举清除。
栅格中圆盘内任一点到最近栅格点的距离不超过 $g/\sqrt2=10.6$ m $<20$ m，
因此\textbf{只要真值落在 $\sigma$ 圆盘内就必然命中}（这由问题一的定位区域理论保证）。
代价最坏约 $25\times3=75$ s，只在最后手段触发。
\end{keybox}

\subsection{效果（四池 120 例 A/B）}

\begin{center}
\begin{tabular}{lcccc}
\toprule
配置 & 最差完成率 & 平均时间/s & 失败算例 & 漏源数 \\ \midrule
本轮改动前 & 0.9954 & 885.3 & 3/120 & 3 \\
\textbf{本轮改动后} & \textbf{0.9978} & \textbf{880.7} & \textbf{1/120} & \textbf{1} \\
\bottomrule
\end{tabular}
\end{center}

\textbf{漏源减少 67\%，时间还略降}——因为这两件工具让机器人不必再反复出车去
"再试一次"，从而省下了无效往返。

\subsection{仅剩的 1 例：几何极限，不是算法缺陷}

\begin{lstlisting}[caption={seed 51029 / ch2 的几何核算}]
source (1767, -219)  r = 1780 m  azimuth = 352.9 deg
radiation direction 348 deg ; dot(dir, outward radial) = 0.996   <-- 几乎完全朝外
lit area inside the arena: 0.11 %  (0.011 km^2 of 10.18 km^2)
deepest lit point on the rim is 27 m beyond the boundary plane
\end{lstlisting}

也就是说：该源的\textbf{可探测区域是一条 27 m 宽的月牙}。
要保证覆盖它，需要在整条 11.3 km 的边界上每 25 m 布一个站位，约 \textbf{450 个站位}；
而当前每例的总时间预算是 28800 s（单源 360000 s），450 站仅移动就需 $450\times120=54000$ s。
\textbf{这在时间预算上不可行}，因此这一例属于问题固有的测度极小的构型，
不应算作算法缺陷。本文据此把\textbf{可达到的完成率上界记作 $1-0.064\%=0.9994$}。

% ===========================================================================
\section{方向二：减少每站检测频道数}

\subsection{实测：检测次数已随频道解算自然下降}

单例剖析（seed 30000，每站的状态分布与本次检测次数）：

\begin{center}
\small
\begin{tabular}{lccccc}
\toprule
站点序 & 未解算 & 已听到 & 已定位 & 已清除 & 本站检测次数 \\ \midrule
1（入场普查） & 15 & 5 & 0 & 0 & 19 \\
2 & 14 & 2 & 4 & 0 & 16 \\
3 & 14 & 0 & 5 & 1 & 14 \\
4 & 12 & 2 & 5 & 1 & 14 \\
5 & 11 & 2 & 6 & 1 & 12 \\
\dots & & & & & \\
中后期各站 & 10 & 0$\sim$2 & 0$\sim$5 & 3$\sim$10 & 10$\sim$12 \\
\bottomrule
\end{tabular}
\end{center}

检测次数由 20 降到 9$\sim$16，说明"已清除 / 已认证"的频道确实被跳过了。

\subsection{新规则"站位附近已全部认证则不再测量"}

逻辑是严格的：对"从未听到"的频道 $c$，在某站位 $p$ 测一次的全部价值，
在于为 $p$ 的认证半径邻域内尚未排除的候选点提供读数；
若该邻域\textbf{全部已排除}，测量不会改变任何结论。判定用认证扫描已维护的增量结构，
代价 $O(\text{邻域点数})\approx800$ 次整数比较，相对一次 6 s 的检测可忽略。

\textbf{实测：几乎不触发}（405 次调用，0 次跳过）。
原因很清楚——认证本身已接近信息下界：每个候选点需要 3$\sim$4 条环绕读数，
而间距 900$\sim$1000 m 的格点恰好只提供这么多。\textbf{再压缩只能削弱认证本身。}

% ===========================================================================
\section{方向三：顺路清除（与巡回规划的混合）}

\subsection{动机}

纯扫描配置的时间构成（16 例，6795 s/例 $=$ 543 s/源）：

\begin{center}
\begin{tabular}{lcccc}
\toprule
任务类 & 次数/例 & 行程/m & 占总行程 & 占总时间 \\ \midrule
普查站位 & 18.6 & 11884 & 53.7\% & 50.5\% \\
\textbf{清除} & 12.7 & \textbf{9060} & \textbf{40.9\%} & 38.0\% \\
补测站位 & 6.8 & 2487 & 11.2\% & 13.6\% \\
入场普查 & 1.0 & 0 & 0\% & 1.8\% \\
\bottomrule
\end{tabular}
\end{center}

清除行程已成为最大的单项（12.7 个源、平均每个 713 m）。
但清除必须走到离源 20 m 以内，"必须走"的部分省不掉，
能省的是\textbf{从扫描路线专门拐出去再拐回来的往返}。

\subsection{实现与实测：负优化}

\begin{center}
\begin{tabular}{lcccc}
\toprule
配置 & 完成率 & 最差 & 平均时间/s & 行程/m \\ \midrule
关闭（基线） & \textbf{1.0000} & \textbf{1.0000} & \textbf{540.3} & 22343 \\
开启，半径 350 m & 0.9972 & 0.9333 & 544.4 & 22741 \\
开启，半径 520 m & 0.9946 & 0.9333 & 573.3 & 24935 \\
开启，半径 750 m & 0.9946 & 0.9333 & 599.3 & 27091 \\
\bottomrule
\end{tabular}
\end{center}

\textbf{完成率与时间同时变差。} 原因是机会式清除打断了扫描，
而 \texttt{clear\_channel} 自身的归航与末段图案带来的额外移动
\textbf{超过了它省下的那趟绕行}；同时推迟了扫描的发现进度
（与上一轮"分阶段调度"失败的机理同源）。

\begin{warnbox}[title=结论]
"顺路清除"在本题\textbf{不成立}。清除不是一个可以"顺手"完成的动作：
它内部含有归航、图案、复测等一串子过程，其代价远大于"走到那里"本身。
这是本项目第二次在"搜索—清除耦合"上得到负结果，两次的机理是一致的。
\end{warnbox}

% ===========================================================================
\section{混合策略的真正答案：认证搜索兼任定向源发现器}

\subsection{发现}

论文配置的 24 个站位里，只有约 7 个来自格点，其余 17 个由认证搜索补出。
把认证搜索整个去掉、改用更密的格点扫描，结果如下（四池各 30 例，按最差池判定）：

\begin{center}
\footnotesize
\begin{tabular}{p{3.6cm}p{1.3cm}p{1.5cm}p{1.5cm}p{1.1cm}p{3.6cm}}
\toprule
配置 & 站位数 & 最差完成率 & 时间/s & 失败 & 各池完成率（A/B/C/D） \\ \midrule
P4 论文（$s{=}1100$ $+$ 认证搜索） & 24.0 & \textbf{0.9978} & 880.7 & \textbf{1/120} & 1.000 / 0.998 / 1.000 / 1.000 \\
$s{=}1000$ 纯扫描 & 12.8 & 0.9833 & \textbf{533.8} & 14/120 & 1.000 / 0.983 / 0.990 / 0.986 \\
$s{=}900$ 纯扫描 & 18.3 & 0.9751 & 554.5 & 17/120 & 0.993 / 0.982 / 0.975 / 1.000 \\
$s{=}950$ 纯扫描 & 12.9 & 0.9803 & 515.1 & 21/120 & 0.987 / 0.980 / 0.986 / 0.986 \\
$s{=}1000$ $+$ 认证搜索（预算 6） & 13.0 & 0.9833 & 539.4 & 14/120 & --- \\
\bottomrule
\end{tabular}
\end{center}

\textbf{关键观察：认证搜索的补站在 $s=1100$ 时有 17 个，在 $s=1000$ 时几乎为 0。}
为什么？因为间距 $\le$ 认证半径 958 m 时，候选点所在格点三角形的三个顶点
都在认证半径内，格点自身就能完成认证；而 $s=1100$ 时不行，于是补站很多。

那为什么 $s=1100+$17 补站反而\textbf{完成率最高}？因为那 17 个补站是
\textbf{被"证书最弱"这一信息驱动}放置的——而"证书最弱"的位置，
恰恰就是\textbf{定向源最可能藏身}的位置。\textbf{认证搜索兼任了定向源发现器。}

\subsection{贴边环：方向正确但不足以补齐}

贴边朝外的定向源，其受光月牙必与半径 1750 m 的圆环相交，
所以在环上均匀补几个站位是"对症"的。实测：

\begin{center}
\begin{tabular}{lccc}
\toprule
配置 & 平均站位数 & 最差完成率 & 平均时间/s \\ \midrule
$s{=}1000$，无贴边环 & 12.8 & 0.9833 & 533.8 \\
$s{=}1000$，贴边环 6 & 18.7 & 0.9879 & 558.3 \\
$s{=}1000$，贴边环 9 & 21.3 & 0.9849 & 561.9 \\
$s{=}1000$，贴边环 12 & 24.2 & 0.9869 & 582.4 \\
$s{=}1000$，贴边环 16 & 28.0 & 0.9905 & 583.2 \\
P4 $+$ 贴边环 6 & --- & 0.9978 & 906.9 \\
\bottomrule
\end{tabular}
\end{center}

贴边环把最差完成率从 0.9833 抬到 0.9905，但\textbf{补不回 P4 的 0.9978}，
而且站位数涨到 28——说明\textbf{均匀补环不如"信息驱动的定向补站"}。

% ===========================================================================
\section{最终前沿：完成度与时间的真实交换}

\begin{center}
\begin{tabular}{lccc}
\toprule
配置 & 最差完成率 & 平均时间/s & 失败算例 \\ \midrule
\textbf{P4 $+$ relocate $+$ exhaustive（推荐）} & \textbf{0.9978} & 880.7 & \textbf{1/120} \\
P4 $+$ 贴边环 6 & 0.9978 & 906.9 & 1/120 \\
$s{=}1000$ $+$ 贴边环 16 & 0.9905 & 583.2 & 12/120 \\
$s{=}1000$ 纯扫描 & 0.9833 & 533.8 & 14/120 \\
$s{=}950$ 纯扫描 & 0.9803 & 515.1 & 21/120 \\
\bottomrule
\end{tabular}
\end{center}

\begin{warnbox}[title=诚实结论]
\textbf{"完成度完整"与"时间大幅压缩"没有同时达成。}
本轮在完成度上取得了实质进步（漏源 3 $\to$ 1，且时间持平），
但把时间从 880 s 压到 530$\sim$580 s 的代价是\textbf{最差完成率下降 0.7$\sim$1.5 个百分点}，
且 120 例中会出现 12$\sim$14 例未能清完。

\textbf{没有找到}既能保持 0.9978 以上、又能把时间压到 600 s 以内的配置。
从数据看，这不是调参不足，而是结构性的：P4 的 17 个认证补站既是
"认证"也是"定向源发现"，去掉它们就同时失去了发现能力。
若要在二者间取得更优的折中，唯一未探索的方向是\textbf{把"发现"与"认证"
明确分成两个目标函数来联合布站}（当前是单一目标：单位代价新认证的候选点数）。
\end{warnbox}

% ===========================================================================
\section{完整尝试清单（本轮）}

\begin{center}
\footnotesize
\scriptsize
\begin{longtable}{p{0.5cm}p{3.4cm}p{3.2cm}p{2.3cm}p{3.1cm}}
\toprule
序号 & 尝试 & 结果 & 判定 & 依据 \\ \midrule
\endhead
R1 & relocate：横向偏移二分 & 单条方位的频道拿到第二条方位 & \textbf{有效} & 4 例已知失败全修复 \\
R2 & relocate 状态机修正 & 候选点从"生成即丢"变为"逐点访问" & \textbf{关键修复} & 单例追踪 \\
R3 & relocate 偏移由小到大 & 局部行程减半但总行程反增、完成率降 & 负优化 & 24 例对照 \\
R4 & relocate 预算（档位数/爬行次数） & 完全无差别 & 无效 & 24 例对照 \\
R5 & exhaustive：末段栅格穷举清除 & 定位到 42 m 却清不掉的源被清掉 & \textbf{有效} & 单例 + 四池 A/B \\
R6 & 硬上限 14 $\to$ 24 & 配合 R5 使用 & 有效 & 同上 \\
R7 & skip\_certified：附近已认证则不测 & 逻辑严格，但几乎不触发 & 无效（近下界） & 405 次调用 0 次跳过 \\
R8 & 顺路清除（半径 350/520/750 m） & 完成率与时间同时变差 & \textbf{负优化} & 24 例对照 \\
R9 & 纯扫描替代认证搜索（$s{=}900/950/1000$） & 时间 $-38\%$，最差完成率 $-1.5$ pp & 权衡，不采用 & 四池 120 例 \\
R10 & 贴边环 6/9/12/16 & 最差 0.9833 $\to$ 0.9905，但补不回 0.9978 & 部分有效 & 四池 120 例 \\
R11 & 认证预算（0/2/4/6/12/40） & 几乎无差别 & 无效 & 多轮对照 \\
R12 & 清除偏置 $\times$ 定位门槛（12 组） & cb500/ls400 最优 & 参数确认 & 24 例 × 12 组 \\
\bottomrule
\end{longtable}
\end{center}

% ===========================================================================
\section{复现方式}

\begin{lstlisting}[language=bash]
cd code
python multi_pool.py 30          # 四池 120 例：候选配置总览（本报告主表）
python ab_relocate.py 30         # 本轮两个新机制的 A/B（漏源 3 -> 1）
python rim_ring_scan.py 30       # 贴边环规模扫描
python relocate_budget.py 24     # relocate 预算
python measure_profile.py 30000  # 每站检测次数剖析
python find_missed.py 40         # 列出所有漏源及其几何特征
\end{lstlisting}

配置：论文默认使用 \texttt{make\_data.P4}（现已包含本轮的两个新机制）；
\texttt{P4\_PERFECT} 与 \texttt{P4\_FAST} 是"省时"预设，其完成率代价见第 6 节，按需选用。

\end{document}
